\documentclass{article}
\usepackage{hyperref}

\begin{document}



\section{Linear Actuator Options}


This document covers the types of linear actuators considered for the front 
oscillation and stabilization of the UAV launch system. Multiple types were considered, 
ranging from prebuilt purchased models to possible self-configurable options. 
The main factors considered are cost, linear force, linear speed, controllability and accuracy.
A rotational force at the base of the launcher was also considered. 

\subsection{Motor With Lead Screw or Ball Screw}


A motor can be connected to a lead screw to transfer rotation motion and force into linear. 
This can be done using differing motors that are explored in detail in section 2.
The main two mechanisms to cause this conversion is a lead screw or a ball screw. 
They both work on similar principles using a motor and gear box to rotate a screw forward and backwards. 

The lead screw uses a threaded shaft and nut mechanism to convert the motion, a ball screw 
instead uses a threaded shaft and ball nut to convert the motion. Ball screws in turn have significantly 
lower friction loss. The advantages and disadvantages of both are shown below. 
\subsubsection{ Lead Screw Advantages}

\begin{itemize}
    \item Low cost.
    \item High mechanical advantage.
    \item High linear force can be achieved.
    \item Simple mechanical construction.
    \item Relatively simple control system.
    \item A lead screw can be self-locking depending on its lead angle.
    \item A self-locking screw can help hold a vertical load without
    continuously powering the motor.
    \item Good for vertical loads to remain stationary. 
    \item Less sensitive to contamination. 
\end{itemize}

\subsubsection{Lead Screw Disadvantages}

\begin{itemize}
    \item Lower efficiency than a ball screw.
    \item Higher friction and heat generation.
    \item Lower maximum speed than many ball-screw systems.
    \item Backlash can reduce positioning accuracy.
    \item Position feedback is required for accurate positioning.
    \item High-speed operation can reduce service life.
    \item Requires more motor torque for similar linear force 
\end{itemize}


\subsubsection{Ball Screw Advantages}

\begin{itemize}
    \item Higher efficiency than a lead screw.
    \item Higher achievable linear speed.
    \item Lower friction.
    \item Lower heat generation.
    \item High force capability.
    \item Good positioning accuracy when combined with position feedback.
    \item Suitable for rapid and repeated motion.
\end{itemize}

\subsubsection{Ball Screw Disadvantages}

\begin{itemize}
    \item Higher cost than a lead screw.
    \item Ball screws can be back-driven by a vertical load.
    \item A brake or continuous motor torque may be required to hold a
    vertical load.
    \item Requires lubrication and maintenance.
    \item More mechanically complex than a lead screw.
    \item Position feedback is required for accurate positioning.
    \item Requires more careful alignment 
\end{itemize}


\section{Motors Paired with Lead/Ball Screw }

While using a ball or lead screw a DC, Stepper or Servo motor can be used.
This document will not be going into detail of the functionality or mechanics or each.
Below is a list of their advantage or disadvantages, when connected to either lead or ball screws. 

\subsection{Dc Motor}
\subsubsection{ DC Advantages}

\begin{itemize}
    \item Relatively low cost.
    \item Simple electrical control.
    \item High rotation low speed.
    \item Direction can be reversed
\end{itemize}

\subsubsection{DC Disadvantages}

\begin{itemize}
    \item Does not inherently know position, requires an encoder
    \item Can overshoot position easily. 
    \item Speed and torque affected by load
    \item Lead-screw friction reduces efficiency.
    \item Maximum linear speed is limited.
    \item Holding position may require continuous motor torque or brake. 
    \item may require gear box to achieve required torque. 
\end{itemize}

\subsection{Stepper Motor}
\subsubsection{Stepper Advantages}

\begin{itemize}
    \item Relatively inexpensive  
    \item Simple position commands 
    \item Does not require encoder 
    \item High holding torque 
    \item Can be very smooth and command resolution 
\end{itemize}


\subsubsection{ Stepper Disadvantages}

\begin{itemize}
    \item Problematic if load exceeds available torque 
    \item Open-loop can not confirm position 
    \item Torque diminishes as rotation speed increases 
    \item Can generate significant heat while stationary 
    \item May require larger motor to achieve high-speed  
    \item Brake or holding mechanism may be necessary for vertical load. 
\end{itemize}

\subsection{Servo Motor}
\subsubsection{Servo Advantages}

\begin{itemize}
    \item Closed-loop position control.
    \item High positioning accuracy and repeatability.
    \item The controller can compensate for changes in load.
    \item Good dynamic performance, rapid acceleration and deceleration. 
    \item  Automatically corrects position errors. 
    \item  Good torque capability over broad operation range. 
\end{itemize}

\subsubsection{Servo Disadvantages}

\begin{itemize}
    \item Higher cost than a basic DC motor system.
    \item Requires a compatible servo drive/controller 
    \item More complex setup. 
    \item Can overshoot or oscillate without proper tuning. 
    \item Holding torque can generate high heat. 
    \item A mechanical brake may still be required to hold load without power. 
\end{itemize}



\section{Linear Motor}

A linear motor produces linear motion directly through electromagnetic
forces rather than converting rotary motion into linear motion through a
screw or other mechanical transmission.

\subsubsection{Advantages}

\begin{itemize}
    \item Direct generation of linear motion.
    \item Very high achievable linear speed.
    \item High acceleration and deceleration.
    \item High positioning accuracy.
    \item High repeatability.
    \item No screw or gearbox backlash.
    \item Low mechanical wear.
    \item Excellent dynamic performance.
\end{itemize}

\subsection{}subsection{Disadvantages}

\begin{itemize}
    \item Significantly higher cost than conventional actuators.
    \item Requires sophisticated control electronics.
    \item Position feedback is typically required.
    \item Does not inherently provide a self-locking mechanism.
    \item A brake or continuous control force may be required for vertical
    loads.
    \item More difficult to integrate mechanically.
    \item Generally unnecessary when extremely high speed is not required.
\end{itemize}

\section{Basic Prebuilt Linear Actuator}

Basic prebuilt electric linear actuator were explored as a possible 
option for the front rail actuation. However this was quickly shown 
to be not feasible due to the slow linear speed of most models making it 
nearly unusable for the required launch time frame. 

\section{Rotation Torque Instead of Linear Force}

In the application of the launch system a motor mounted at the point axis 
of rotation is also a feasible solution. This inherently changes some of the 
required specification need from the motors. Instead of linear speed it would be 
angular speed as well as applied torque. The general torque would need to be significantly 
greater due to the force applied at the base. 

This option is still in considerations as it the motor may be larger and more expensive 
but could be compensated with build gear box and no longer the need for a lead or ball screw 
mechanism


\section{Summary}

In summary, there are multiple options for actuation of the front end of the UAV launch 
system. The main leading option being a ball screw with a servo motor. 
\newline
\newline 
The lead screw is likely 
to cause complication with friction at high speeds. The Dc and Stepper motor although cheaper 
will likely lead to lower position and control accuracy then needed for this application. 
The linear motor although the best for force and controllability is largely too 
expensive of an option for the budget of this project. The basic linear actuator 
is incapable of getting to required speed to control the actuation properly.
The base rotation with a motor and gearbox is a viable solution. \newline

Final part selection Will be completed for both the full size UAV launcher and scaled down
version in later documentation, using this documents information along with calculations taken from
"UAV Launch Actuator Values". 

\begin{thebibliography}{9}

\bibitem{isel}
isel USA, ``The Differences Between Lead and Ball Screw Actuators.''
[Online]. Available: \url{https://www.isel-us.com/blog/the-differences-between-lead-and-ball-screw-actuators}.

\bibitem{amc}
Advanced Motion Controls, ``DC Servo Motor.''
[Online]. Available: \url{https://www.a-m-c.com/dc-servo-motor/}.

\bibitem{pihut}
The Pi Hut, ``What's the Difference Between DC, Servo and Stepper Motors.''
[Online]. Available: \url{https://thepihut.com/blogs/raspberry-pi-tutorials/whats-the-difference-between-dc-servo-amp-stepper-motors}.

\bibitem{tecnotion}
Tecnotion, ``What Is a Linear Motor?''
[Online]. Available: \url{https://www.tecnotion.com/faq/what-is-a-linear-motor/}.

\bibitem{acv}
ACV Group, ``Understanding How Electric Actuators Work.''
[Online]. Available: \url{https://acv-group.com/understanding-how-electric-actuators-work/}.

\end{thebibliography}

\end{document}
